% GENERATED FILE. Edit canonical lessons, then run npm run generate:books.
% canonical-source-hash: fnv1a64:2c40c90d0ed58c40
% canonical-lessons: GE-C403-anspringen, GE-C403-batterie, GE-C403-starthilfe, GE-C403-abschleppwagen, GE-C403-mechaniker

\chapter{Der Motor springt nicht an}
\label{ch:ge-der-motor-springt-nicht-an}
\hlchaptermodality{403}

\begin{chapteropening}
\textbf{By the end of this chapter:} \emph{I can talk about a car that will not start: starting, the battery, a jump start, the tow truck and the mechanic.}

Complete the last lesson of chapter 403: I can talk about a car that will not start: starting, the battery, a jump start, the tow truck and the mechanic.
\end{chapteropening}

\section[anspringen]{anspringen — to start (of an engine)}

\label{lesson:GE-C403-anspringen}



\begin{quote}
Before the new one: say the German for the petrol pump, then the German for the oil change.
\end{quote}

\subsection*{You'll want to know: anspringen}
\textbf{anspringen} — \textquotedblleft{}to start (of an engine)\textquotedblright{}.

\begin{quote}
Im Winter springt unser Auto morgens schlecht an.
\end{quote}

\begin{quote}
\emph{In winter our car is hard to start in the mornings.}
\end{quote}

Separable; Perfekt with \emph{sein}: \emph{der Motor ist nicht angesprungen}.

\begin{cousinweb}
\emph{anspringen} is built from \emph{an} 'on' and \emph{springen} 'to jump', inherited from Proto-Germanic \emph{springaną} (EtymDB): an engine that jumps into life. English \emph{spring} is a cousin.
\end{cousinweb}

\begin{grammarlens}[title={What you've built}]
One more word for telling, explaining and describing at B1.
\end{grammarlens}

\subsection*{Your turn}
\begin{itemize}
\raggedright
  \item \emph{Say it:} \emph{anspringen}
  \item \emph{Say it:} \emph{anspringen}, once more
  \item \emph{Say it:} a sentence of your own with \emph{anspringen}
  \item \emph{Recall:} say the German for the petrol pump, then the German for the oil change, then say \emph{anspringen} again
\end{itemize}

\begin{tcolorbox}[breakable,colback=teal!4,colframe=teal!35!black,title={Before you move on}]
What does anspringen mean? (\textquotedblleft{}To start (of an engine)\textquotedblright{}.) Say it once more.
\end{tcolorbox}
\section[die Batterie]{die Batterie — the battery}

\label{lesson:GE-C403-batterie}



\begin{quote}
Before the new one: say the German for the oil change, then the German for to start.
\end{quote}

\subsection*{You'll want to know: die Batterie}
\textbf{die Batterie} — \textquotedblleft{}the battery\textquotedblright{}.

\begin{quote}
Die Batterie ist leer, weil das Licht die ganze Nacht an war.
\end{quote}

\begin{quote}
\emph{The battery is flat because the light was on all night.}
\end{quote}

Plural: \emph{die Batterien}. The battery in a phone is usually \emph{der Akku}.

\begin{cousinweb}
\emph{Batterie} was borrowed from French \emph{batterie}, first 'a beating', then a row of guns, from \emph{battre} 'to beat' (EtymDB), which goes back to Latin \emph{battuere} 'to beat' (Lewis \& Short).
\end{cousinweb}

\begin{grammarlens}[title={What you've built}]
One more word for telling, explaining and describing at B1.
\end{grammarlens}

\subsection*{Your turn}
\begin{itemize}
\raggedright
  \item \emph{Say it:} \emph{die Batterie}
  \item \emph{Say it:} \emph{die Batterie}, once more
  \item \emph{Say it:} a sentence of your own with \emph{die Batterie}
  \item \emph{Recall:} say the German for the oil change, then the German for to start, then say \emph{die Batterie} again
\end{itemize}

\begin{tcolorbox}[breakable,colback=teal!4,colframe=teal!35!black,title={Before you move on}]
What does die Batterie mean? (\textquotedblleft{}The battery\textquotedblright{}.) Say it once more.
\end{tcolorbox}
\section[die Starthilfe]{die Starthilfe — the jump start}

\label{lesson:GE-C403-starthilfe}



\begin{quote}
Before the new one: say the German for to start, then the German for the battery.
\end{quote}

\subsection*{You'll want to know: die Starthilfe}
\textbf{die Starthilfe} — \textquotedblleft{}the jump start\textquotedblright{}.

\begin{quote}
Ein Nachbar gibt uns mit seinem Auto Starthilfe.
\end{quote}

\begin{quote}
\emph{A neighbour gives us a jump start with his car.}
\end{quote}

No plural in this sense. \emph{Das Starthilfekabel} — the jump lead.

\begin{cousinweb}
\emph{Starthilfe} is built from \emph{Start}, usually traced to English \emph{start} (EtymDB lists \emph{Start} without a source), and \emph{Hilfe} 'help', from the inherited verb \emph{helfen}.
\end{cousinweb}

\begin{grammarlens}[title={What you've built}]
One more word for telling, explaining and describing at B1.
\end{grammarlens}

\subsection*{Your turn}
\begin{itemize}
\raggedright
  \item \emph{Say it:} \emph{die Starthilfe}
  \item \emph{Say it:} \emph{die Starthilfe}, once more
  \item \emph{Say it:} a sentence of your own with \emph{die Starthilfe}
  \item \emph{Recall:} say the German for to start, then the German for the battery, then say \emph{die Starthilfe} again
\end{itemize}

\begin{tcolorbox}[breakable,colback=teal!4,colframe=teal!35!black,title={Before you move on}]
What does die Starthilfe mean? (\textquotedblleft{}The jump start\textquotedblright{}.) Say it once more.
\end{tcolorbox}
\section[der Abschleppwagen]{der Abschleppwagen — the tow truck, the breakdown lorry}

\label{lesson:GE-C403-abschleppwagen}



\begin{quote}
Before the new one: say the German for the battery, then the German for the jump start.
\end{quote}

\subsection*{You'll want to know: der Abschleppwagen}
\textbf{der Abschleppwagen} — \textquotedblleft{}the tow truck, the breakdown lorry\textquotedblright{}.

\begin{quote}
Der Abschleppwagen bringt das kaputte Auto in die Werkstatt.
\end{quote}

\begin{quote}
\emph{The tow truck takes the broken-down car to the garage.}
\end{quote}

Same form in the plural.

\begin{cousinweb}
\emph{Abschleppwagen} is built from \emph{abschleppen} 'to tow away', from \emph{schleppen} 'to drag', which came in from Middle Low German \emph{slēpen} (EtymDB), and \emph{Wagen} 'vehicle', inherited from Proto-Germanic \emph{wagnaz} (EtymDB).
\end{cousinweb}

\begin{grammarlens}[title={What you've built}]
One more word for telling, explaining and describing at B1.
\end{grammarlens}

\subsection*{Your turn}
\begin{itemize}
\raggedright
  \item \emph{Say it:} \emph{der Abschleppwagen}
  \item \emph{Say it:} \emph{der Abschleppwagen}, once more
  \item \emph{Say it:} a sentence of your own with \emph{der Abschleppwagen}
  \item \emph{Recall:} say the German for the battery, then the German for the jump start, then say \emph{der Abschleppwagen} again
\end{itemize}

\begin{tcolorbox}[breakable,colback=teal!4,colframe=teal!35!black,title={Before you move on}]
What does der Abschleppwagen mean? (\textquotedblleft{}The tow truck, the breakdown lorry\textquotedblright{}.) Say it once more.
\end{tcolorbox}
\section[der Mechaniker]{der Mechaniker — the mechanic}

\label{lesson:GE-C403-mechaniker}



\begin{quote}
Before the new one: say the German for the jump start, then the German for the tow truck.
\end{quote}

\subsection*{You'll want to know: der Mechaniker}
\textbf{der Mechaniker} — \textquotedblleft{}the mechanic\textquotedblright{}.

\begin{quote}
Der Mechaniker findet den Fehler nach zehn Minuten.
\end{quote}

\begin{quote}
\emph{The mechanic finds the fault after ten minutes.}
\end{quote}

Same form in the plural. A woman: \emph{die Mechanikerin}.

\begin{cousinweb}
\emph{Mechaniker} is a learned word built on New Latin \emph{mechanicus}, from Greek \emph{mēkhanikós} 'of machines' (EtymDB), from \emph{mēkhan\'{\={e}}} 'device'.
\end{cousinweb}

\begin{grammarlens}[title={What you've built}]
One more word for telling, explaining and describing at B1.
\end{grammarlens}

\subsection*{Your turn}
\begin{itemize}
\raggedright
  \item \emph{Say it:} \emph{der Mechaniker}
  \item \emph{Say it:} \emph{der Mechaniker}, once more
  \item \emph{Say it:} a sentence of your own with \emph{der Mechaniker}
  \item \emph{Recall:} say the German for the jump start, then the German for the tow truck, then say \emph{der Mechaniker} again
\end{itemize}

\begin{tcolorbox}[breakable,colback=teal!4,colframe=teal!35!black,title={Before you move on}]
What does der Mechaniker mean? (\textquotedblleft{}The mechanic\textquotedblright{}.) Say it once more.
\end{tcolorbox}
